\chapter*{Vorwort: Was ein Modell kann}
\addcontentsline{toc}{chapter}{Vorwort}
\markboth{Vorwort}{Vorwort}

Die Biomathematik ist das Gebiet, in dem die Mathematik am sichtbarsten \kursiv{scheitert} --- und gerade deshalb am meisten leistet.

\begin{merksatz}
\boxtitel{Der Ausgangspunkt}
\begin{center}
\begin{minipage}{0.9\textwidth}
Jedes Modell dieses Bandes ist falsch. Kein Bakterium teilt sich exakt nach einer Exponentialfunktion, keine Epidemie verläuft wie eine Gleichung, kein Wirkstoff verteilt sich in drei wohldefinierten Behältern.
\end{minipage}
\end{center}

\begin{center}
\begin{minipage}{0.88\textwidth}
\kursiv{Die Frage ist daher nie, ob ein Modell stimmt --- es stimmt nie. Die Frage lautet: Welche Auskunft gibt es trotzdem, und ab wann führt es in die Irre? Das ist die Leitfrage dieses Bandes, und sie steht in jedem Kapitel als eigener Abschnitt.}
\end{minipage}
\end{center}
\end{merksatz}

\section*{Warum sich das lohnt}

\begin{merksatz}
\boxtitel{Vier Dinge, die ein Modell leistet}
\begin{center}
{\small
\renewcommand{\arraystretch}{1.55}
\begin{tabular}{@{}P{3.0cm}P{5.6cm}@{}}
\toprule
\textbf{Zusammen\-hänge sichtbar machen} & Warum wirkt eine Impfquote von $80\,\%$ anders als eine von $60\,\%$? Die Antwort ist eine Schwelle, und man sieht sie nur in der Gleichung.\\
\textbf{Größen abschätzen} & Wie lange dauert es, bis ein Wirkstoffspiegel sein Gleichgewicht erreicht? Vier Halbwertszeiten --- das folgt aus der Rechnung, nicht aus der Erfahrung.\\
\textbf{Annahmen prüfen} & Ein Modell zwingt dazu, alle Annahmen aufzuschreiben. Widersprüche werden dabei sichtbar.\\
\textbf{Überraschungen erzeugen} & Einige Ergebnisse dieses Bandes widersprechen der Anschauung --- und sind trotzdem richtig.\\
\bottomrule
\end{tabular}}
\end{center}

\begin{center}
\begin{minipage}{0.88\textwidth}
\kursiv{Zur letzten Zeile ein Vorgriff: Eine einfache Gleichung für eine Insektenpopulation --- eine Zeile, drei Zeichen --- erzeugt bei bestimmten Parameterwerten einen Verlauf, der sich nie wiederholt und von keiner Rechnung vorhersagen lässt. Kapitel~3 führt das vor. Es war eine der folgenreichsten Entdeckungen der Mathematik des zwanzigsten Jahrhunderts, und sie stammt aus der Biologie.}
\end{minipage}
\end{center}
\end{merksatz}

\begin{warnung}
\boxtitelr{Der Aufbau jedes Kapitels}
\begin{center}
\renewcommand{\arraystretch}{1.5}
\begin{tabular}{@{}P{2.4cm}P{6.2cm}@{}}
\toprule
\textbf{Die Annahme} & Was genau wird unterstellt? In Worten, vor jeder Formel.\\
\textbf{Die Gleichung} & Aus der Annahme aufgestellt, Schritt für Schritt.\\
\textbf{Die Lösung} & Vollständig hergeleitet, nicht nachgeschlagen.\\
\textbf{Die Deutung} & Was bedeuten die Parameter in der Sache?\\
\textbf{Der Bruchfall} & Wo versagt das Modell, und woran erkennt man es?\\
\bottomrule
\end{tabular}
\end{center}

\begin{center}
\begin{minipage}{0.88\textwidth}
\kursiv{Der letzte Punkt ist derselbe wie in Band 17. Ein Modell, dessen Grenzen man nicht kennt, wendet man blind an --- und in der Biomathematik sind diese Grenzen oft näher, als die Eleganz der Gleichungen vermuten lässt.}
\end{minipage}
\end{center}
\end{warnung}

\section*{Was Sie mitbringen müssen}

\begin{merksatz}
\boxtitel{Die Voraussetzungen}
\begin{center}
{\small
\renewcommand{\arraystretch}{1.5}
\begin{tabular}{@{}P{3.0cm}P{5.6cm}@{}}
\toprule
\textbf{Ableiten} & Was $f'(x)$ bedeutet und wie man Polynome, Exponentialfunktionen und Produkte ableitet. Band 3.\\
\textbf{Integrieren} & Stammfunktionen der Grundfunktionen; Trennung der Veränderlichen wird in Kapitel~1 vorgeführt.\\
\textbf{Exponential und Logarithmus} & Die Rechenregeln. Band 0, Teil~IV.\\
\textbf{Matrizen} & Erst ab Kapitel~6, und dort nur $2\times2$. Band 2.\\
\textbf{Statistik} & Nicht nötig. Wo Daten vorkommen, wird auf Band 17 verwiesen.\\
\bottomrule
\end{tabular}}
\end{center}

\begin{center}
\begin{minipage}{0.88\textwidth}
\kursiv{Band 14 \enquote{Die Gesetze der Veränderung} behandelt Differentialgleichungen allgemein und systematisch. Dieser Band setzt ihn nicht voraus: Jede benötigte Lösungsmethode wird an der Stelle hergeleitet, an der sie gebraucht wird. Wer die Theorie dahinter sehen will, findet sie dort.}
\end{minipage}
\end{center}
\end{merksatz}

\section*{Neustadt und sein Klinikum}

\begin{merksatz}
\boxtitel{Der Leitfall}
Wie in Band 17 läuft ein einziger erfundener Fall durch alle fünfzehn Kapitel.

\begin{center}
\begin{minipage}{0.9\textwidth}
In der Stadt \textbf{Neustadt} mit $100\,000$ Einwohnern breitet sich ein bakterieller Erreger aus. Am Klinikum wird er in Kulturen untersucht, mit einem Antibiotikum behandelt --- und nach einigen Monaten treten erste unempfindliche Stämme auf.
\end{minipage}
\end{center}

\begin{center}
{\small
\renewcommand{\arraystretch}{1.5}
\begin{tabular}{@{}P{1.6cm}P{2.8cm}P{3.6cm}@{}}
\toprule
\textbf{Teil} & \textbf{Thema} & \textbf{Der Leitfall zeigt}\\
\midrule
I & Wachstum & die Kultur im Labor\\
II & Wechselwirkung & zwei Stämme im selben Nährmedium\\
III & Ausbreitung & die Epidemie in Neustadt\\
IV & Der Körper & das Antibiotikum im Patienten\\
V & Krankheit & die Entstehung der Resistenz\\
VI & Stand & was die Forschung heute kann\\
\bottomrule
\end{tabular}}
\end{center}

\begin{center}
\begin{minipage}{0.88\textwidth}
\kursiv{Die Zahlen sind erfunden und so gewählt, dass sie sich von Hand rechnen lassen. Sie sind aber der Größenordnung nach realistisch --- Verdopplungszeiten, Halbwertszeiten und Reproduktionszahlen liegen in Bereichen, die man aus der Literatur kennt.}
\end{minipage}
\end{center}
\end{merksatz}

\section*{Drei Warnungen}

\begin{warnung}
\boxtitelr{Erstens: Parameter sind keine Naturkonstanten}
In jeder Gleichung dieses Bandes stehen Buchstaben wie $r$, $K$, $\beta$. Sie sehen aus wie feste Größen. Sie sind es nicht.

\begin{center}
\begin{minipage}{0.9\textwidth}
Eine Wachstumsrate hängt von Temperatur, Nährstoffen und Stamm ab. Eine Übertragungsrate hängt vom Verhalten der Menschen ab --- und ändert sich, sobald über die Epidemie berichtet wird.
\end{minipage}
\end{center}

\kursiv{Wer einen Parameter aus Daten schätzt, schätzt ihn für die Umstände, unter denen die Daten entstanden.}
\end{warnung}

\begin{warnung}
\boxtitelr{Zweitens: Genauigkeit ist nicht Richtigkeit}
Ein Modell liefert Zahlen mit beliebig vielen Stellen. Das verführt dazu, sie für genau zu halten.

Bei Epidemiemodellen ist der Fehler besonders groß: Kleine Änderungen der Übertragungsrate verändern den Verlauf erheblich, weil sie sich über viele Generationen aufschaukeln. Kapitel~7 rechnet das vor.
\end{warnung}

\begin{warnung}
\boxtitelr{Drittens: Das einfachste Modell ist oft das nützlichste}
Man kann jedes Modell verfeinern --- Altersgruppen, räumliche Verteilung, Verhaltensänderung. Jede Verfeinerung fügt Parameter hinzu, die man nicht kennt.

\begin{center}
\begin{minipage}{0.88\textwidth}
\kursiv{Ein Modell mit drei Gleichungen, dessen Verhalten man versteht, ist einem mit dreißig überlegen, dessen Ergebnis man nur glauben kann. Die Aufgabe der Modellierung ist nicht, die Wirklichkeit nachzubauen --- dafür gibt es die Wirklichkeit.}
\end{minipage}
\end{center}
\end{warnung}

\section*{Der Aufbau}

\begin{merksatz}
\boxtitel{Sechs Teile, fünfzehn Kapitel}
\begin{center}
{\small
\renewcommand{\arraystretch}{1.45}
\begin{tabular}{@{}P{1.0cm}P{3.0cm}P{4.6cm}@{}}
\toprule
\textbf{Teil} & \textbf{Thema} & \textbf{Die Leitfrage}\\
\midrule
I & Wachstum (1--3) & Wie viel wird daraus?\\
II & Wechselwirkung (4--6) & Was geschieht, wenn zwei sich beeinflussen?\\
III & Ausbreitung (7--9) & Wann bricht eine Epidemie aus?\\
IV & Der Körper (10--12) & Wie verteilt und wirkt ein Stoff?\\
V & Krankheit (13--14) & Wie wachsen Tumoren, wie entsteht Resistenz?\\
VI & Stand (15) & Was können Modelle heute?\\
\bottomrule
\end{tabular}}
\end{center}
\end{merksatz}

\vspace{10pt}
\begin{flushright}
\itshape\color{grau}
Alle Modelle sind falsch.\\
Einige sind brauchbar.\\[8pt]
\normalfont\color{grau}\upshape --- George Box, 1976
\end{flushright}
